% The developed commentary of the concise study, from physical page 5. The
% three interpretations of the expansive study are set beside one another at
% the texts where their answers meet or differ, not summarized one after
% another. Every author is credited only with what he says of his own passage.

\section{The Propers: Three Interpretations Compared}
\zlabel{triptych:concise:commentary:start}

Each of the three interpretations takes one of those questions. The first,
\emph{Rend the Heavens and Come Down}, holds that the prophet's cry is a
prayer God answered in the Incarnation and that the Church still prays it
toward Christ's return. The second, \emph{The Watch at the Door}, holds that
Christ hides the hour so that his servants' readiness never lapses, and that
his command reaches everyone because each one's death brings the end to him.
The third, \emph{The Clay and the Faithful Potter}, holds that a people which
has strayed rests its hope of being found blameless on who God is, Father and
maker and faithful to his call, and not on its own justice. The three are
compatible. They differ in which words of the Mass they press and in what
those words are taken to show.

\subsection{``Thou didst come down'': a prayer granted, or a memory?}

At the cry Jerome explains the wish by what the speakers had been
given until then, which was only promises: ``O that thou wouldst rend the
heavens and come down, and that thou, who art always being promised, wouldst
at last fulfil what was promised.'' The prayer was said, he writes, before
the Saviour came and took our nature from the Virgin. Two verses on he read in
the Hebrew what the Latin of the reading still has, \latin{Descendisti}, and
took the change from wish to statement as the sign of a prayer granted:
\latin{Pulchreque juxta Hebraicum, quasi exauditi essent}, ``beautifully,
according to the Hebrew, as if they had been heard''. What came down is the
Word made flesh, Emmanuel (\work{Commentarii in Isaiam} XVII, PL~24,
cols.~645--646).

Thomas Aquinas, commenting on the same verses, keeps Jerome's answer as one
sense and prefers another. He takes \latin{dirumperes}, ``that thou wouldst
rend'', either as a figure addressed to one who dwells above the heavens or as
``that thou wouldst empty thyself, setting thy majesty aside, and take
flesh''. Of \latin{Descendisti} he gives three senses: the coming
down to judge, the past tense put for the future; the coming down into the
Virgin's womb; or, \latin{melius, secundum litteralem sensum}, ``better,
according to the literal sense'', a proof drawn from the past: if you came
down now, these things would happen as they happened when you came down to
free your people from Egypt (\work{Expositio super Isaiam} c.~64). For Jerome
the past tense is the answer to the prayer. For Aquinas, in the sense he
prefers, it is the people's memory of the Exodus, the ground on which they
dare to ask for it again. Here the two commentators disagree over the literal
sense; because Aquinas keeps the Incarnation among the senses of the verb,
Jerome's answer stands, though not as the literal sense.

The third interpretation hears the same chapter as a confession of sin and an
appeal to God as Father. The Latin of the reading has both
\latin{Descendisti} and the confession that follows it. The English proclaimed
in the United States has the cry and the confession and lacks
\latin{Descendisti}, so
that an assembly hearing it in English hears a people asking God to come, not
saying that he came; on that text the third interpretation stands closer than
the first.

Augustine, at the verse sung before the Gospel, says why the Church still
asks for a Christ already given. The salvation asked for is Christ: \latin{salutare tuum, Christum tuum}, ``Thy salvation, that is,
Thy Christ.'' Then: ``He has already given us his Christ; yet let us still say
to him, Give us thy Christ,'' as we ask for daily bread, and our bread is the
one who said, ``I am the living bread which came down from heaven.'' And a
second reason: \latin{Dedit enim nobis Christum, sed hominem: quem nobis dedit
hominem, eum ipsum nobis daturus est Deum}, ``For he gave us Christ, but as
man; the one he gave us as man, the same he will give us as God.'' No man
could receive him as God; we are sons already in hope, not yet in fact, and
ask to see him ``as the Angels see Him'' (\work{Enarrationes in Psalmos} 84.9).
Cyril of Jerusalem, instructing candidates for baptism, teaches two comings: ``We preach not
one advent only of Christ, but a second also, far more glorious than the
former''; the first ``unobserved, like rain on a fleece'', the second ``His open
coming, which is to be'' (\work{Catechetical Lectures} 15.1).

Cyril's distinction suggests a further step, which is not his. Nothing at
Bethlehem tore the sky or melted the mountains, so the prophet's image fits
the open coming better than the hidden one, while for Jerome the hidden coming
already answered the cry. The two are compatible: what the prayer asked to
see, the sky torn open, is still to come. The Preface gives thanks for both,
the coming \latin{in humilitate carnis assumptae}, in the lowliness of the
flesh he took, and the coming \latin{in suae gloria maiestatis}, in the glory
of his majesty, and the solemn blessing names them together in its
last invocation.

\subsection{``Why hast thou made us to err?'': the ground of the people's hope}

The question that opens the reading sounds like an accusation:
\latin{Quare errare nos fecisti, Domine, de viis tuis; indurasti cor
nostrum?}, ``Why hast thou made us to err, O Lord, from thy ways: why hast
thou hardened our heart?'' (63:17). Jerome and Aquinas both refuse to read it
as saying that God causes sin, and they give different reasons. Jerome says
that God's patience only seems to be the cause: he waits for the sinner's
salvation and does not correct him at once, and the sinner, left uncorrected,
goes further. He compares a physician who hides the knife from his patient
until the moment of cutting; the delay is for the patient's sake, though it
looks like neglect (cols.~643--644). Aquinas puts it briefly: God made
them err \latin{permittendo}, ``by permitting it'', not sending grace and not
correcting; and ``thou hast hardened'' means ``thou hast permitted'' (c.~63).
The English proclaimed in the United States asks why God \emph{lets} his
people wander, close to Aquinas's \latin{permittendo}. Jerome and Aquinas were
explaining the harder Latin.

In the confession the anger stands first: \latin{Ecce tu iratus es, et
peccavimus}, ``behold thou art angry, and we have sinned'' (64:4; Douay 64:5).
The two commentators give the same answer in nearly the same words. Jerome:
``The order is inverted. It is not because thou wast angry that we sinned, but
because we sinned that thou wast angry'' (col.~647). Aquinas: ``an inverted
order: thou wast angry because we sinned: and yet we shall be saved, by thy
mercy'' (c.~64). Both read the end of the verse, ``and we shall be saved'', as
the people's hope, and both place it in mercy; Jerome adds that those who were
always in sin will be saved by mercy alone. At the last verse, where the
people call themselves clay, Jerome turns from their merits to God's
clemency: if we consider our merits, we must despair; if we consider thy
clemency, we dare to pray, because thou art our Father. The pot may not answer
back to the potter, but knowing that they are sons the people dare to say the
next verse, ``Be not very angry, O Lord'' (cols.~648--650).

Set side by side, Aquinas's preferred sense of \latin{Descendisti} and
Jerome's comment on the clay name two grounds on which the people ask: what
God once did for them, and what God is to them. The third interpretation rests
on the second.

The second interpretation hears in the same confession the sleep its Gospel
warns against: ``There is none that calleth upon thy name: that riseth up, and
taketh hold of thee'' (64:6; Douay 64:7), where the Latin verb,
\latin{consurgat}, is to rise up.

\subsection{``Convert us'': who is asked to turn?}

The confession goes on, ``thou hast hid thy face from us''. Aquinas reads the
hidden face as mercy withheld, and Jerome, on the same words, quotes the psalm
sung after the reading: ``shew us thy face, and we shall be saved''
(col.~649). The petition in Isaiah and the psalm's refrain still ask different
things. In the
Latin, Isaiah's people ask God himself to turn back: \latin{Convertere propter
servos tuos}, ``return for the sake of thy servants'' (63:17). The psalm's
second strophe uses the same imperative, \latin{Deus virtutum, convertere},
``Turn again, O God of hosts'' (v.~15), though the American psalm renders it
with an adverb. The refrain asks God to turn \emph{us}: \latin{converte nos}.
Jerome reads Isaiah's petition both ways, \latin{Convertere, ait, Domine, sive
converte nos}, ``Turn, he says, O Lord, or, turn us'', and develops the second
with a clause from the Greek version: having a little of the holy mountain, we
may come to the whole and behold the glory of the Only-begotten
(cols.~643--644). ``Turn us'' is how Jerome takes the verse, not its wording
in any text he translated. Aquinas takes \latin{convertere} as a plea that God
turn from anger to mercy, and sets beside it a psalm that asks God to return
(Ps 90 (89):13) and the close of Lamentations, ``convert us, O Lord, and we
shall be converted'' (Lam 5:21).

The third interpretation rests on the refrain's petition. Augustine's
comment on it is one sentence: ``For averse we have been from Thee, and except
Thou convert us, we shall not be converted'' (79.3). Robert Bellarmine
explains why the refrain asks two things: God does not look on us because we
turn away, and ``as we cannot turn to him without his aid, the prophet asks
for both'', that God turn us and show us his face (\work{Commentary on the
Book of Psalms}, Ps 79, expl.~3). On the psalm of the Alleluia verse Augustine
says it again, at the words ``Thou shalt turn us again, and make us alive'':
``not only our being made alive is from Thee, but our very conversion''
(84.8).

The first interpretation hears in the same psalm the coming one. At the son
of man whom God has confirmed for himself Augustine breaks off his
explanation, ``What
can be more evident?'', and names him with Paul as the one foundation, Christ
Jesus (79.9). Bellarmine reports the same as common ground: ``All agree that
this passage refers literally to Christ'', because the vine was never made
perfect until he came (Ps 79, expl.~15). The second interpretation hears the
people's vow at the end, ``we depart not from thee \ldots\ and we will call
upon thy name'', the reverse of Isaiah's ``none that calleth upon thy name''.
Augustine explains the words between, ``Thou wilt quicken us'': ``for dead we
were, when to earthly things we did cleave''; and to call on the name is to
love: ``Thee we will love'' (79.11).

\subsection{The hidden hour and the command to all}

The Gospel raises two difficulties. Jesus has just said that no one knows the
day, ``neither the angels in heaven, nor the Son'' (Mk 13:32): why should the
hour be hidden at all? And he extends his command to all, though nearly all
who hear it will die before the end. The answers to the first come from
commentaries on the parallel in Matthew. Jerome meets the saying first as the
Arians' proof that the Son was ignorant, and answers that the Word through
whom all times were made cannot be ignorant of one day in them. He knows, but
it does not profit the apostles to know, ``so that, always uncertain of the
Judge's coming, they may live each day as if they were to be judged the
next''. At the command to watch Jerome notices the pronoun: ``He did not say,
`we do not know at what hour the Lord will come', but `you do not know''', and
the uncertainty is kept so that ``they may always believe that he will come''
(\work{Commentarii in Matthaeum} IV, on Mt 24:36 and 24:42--43). Bede, on
Mark 13:32--33, copied both comments nearly word for word. John Chrysostom,
preaching on Matthew, makes the same observation in Greek, ``He said not, `I
know not,' but, `ye know not''', and widens it: ``if the common sort of men
knew when they were to die, they would surely strive earnestly at that hour
\ldots\ Wherefore He made the end of each man's life also uncertain''
(\work{Homilies on Matthew} 77.2--3).

Augustine meets the second difficulty in a letter to Hesychius, bishop of
Salona. ``What I say to you, I say to all'' was said, he writes, to the elect
of Christ's body, from the first hearers to those who will be alive at his
last coming; the day will not find all of them alive, and yet the command
binds them all, because \latin{in quo enim quemque invenerit suus novissimus
dies, in hoc eum comprehendet mundi novissimus dies}, ``in whatever state each
one's own last day finds him, in that state the last day of the world will
take him'' (\work{Epistula} 199.1.2--3). In the same letter he separates
ignorance of the date, which is no fault, from the bad servant's ``My lord is
long a coming'' (Mt 24:48), which is.

Gregory the Great preached on Luke's servants who wait, lamps lit, for their
master to come back from a wedding (Lk 12:35--40). The master's knock, he
says, is sickness: ``The Lord comes when he hastens to judgment; he knocks
when, through the troubles of sickness, he signals that death is near''; and
the servant who opens at once is the one who waits gladly for the Judge. He
then says what keeping awake is: he watches ``who holds the eyes of his mind
open to the sight of the true light'', ``who keeps by his works what he
believes'', ``who drives from himself the darkness of sloth and negligence''
(\work{Homiliae in Evangelia} 13.3). Luke's master may come in the second
watch or the third, and Gregory reads the watches as the ages of a life,
childhood, youth and old age: one who did not keep awake in the first should
keep the second, and one who failed in the second should not neglect the
third, for ``even so, there is no cause for despair'' (13.5). Mark's four watches of the night are a different list.

On Mark's parable itself Bede speaks in his own words. The man gone abroad is
Christ ascended, who left the Church in bodily presence ``yet never left it
without the protection of his divine presence''; the identification he takes
from Gregory's homily on the man of Matthew~25 who travels into a far country
(\work{Homiliae in Evangelia} 9).
The authority given to the servants is the grace of the Holy Spirit, given so
that they may serve in good works; the doorkeeper is the order of pastors.
And on the last verse, ``Not only the apostles and their successors \ldots\ but
all of us are commanded to watch'': ``All of us are bidden to guard earnestly
the doors of our hearts, lest the old enemy break in by suggesting evil'',
``for each one will give account to God for himself'' (\work{In Marci
evangelium expositio} IV, on 13:34--37).

Augustine and Bede put the weight in different places. Augustine moves the
command from the end of history into each life, so that the watch is kept
against one's own death; Bede keeps it at the Lord's return and divides it by
office, the pastors at the door of the Church and every person at the door of
his heart. The objection to Augustine is that the parable speaks of the
master's return, not of each servant's death. His reply is in the sentence
quoted: he keeps the last day of the world and says that it reaches each
person in his own last day, since each will be judged as he is then found.

The other two interpretations use the parable differently. In the first,
the master has already been in his house and has gone abroad: Isaiah's people
ask God to come down; the disciples are told to stay awake for one who has
come and will come back. None of these Fathers sets Isaiah's cry beside
Mark's parable; the Mass reads the two on the same day. The
third draws on Bede's comment on the servants' authority. Taken alone, the
second interpretation could sound as though the servants secured their own
readiness by effort; against that, the third points out that what the master
hands them is
grace for good works, that the will to go out to meet Christ is itself asked
for in the Collect, and that it is Christ who ``will confirm you unto the
end''. The second in turn keeps the third from passivity, since the grace is
given as authority over every work.

\subsection{Blameless on the day: Paul's waiting church}

All three interpretations read the thanksgiving of First Corinthians, and the
same two commentators, John Chrysostom and Thomas Aquinas, appear in all
three. In one
sentence Paul thanks God for gifts already given in Christ and describes the
Corinthians as ``waiting for the manifestation of our Lord Jesus Christ''
(1:7). Chrysostom, in his second homily on the letter, first asks how Paul
can praise a church he is about to call carnal. Either they grew remiss after
receiving the gifts, he answers, or Paul speaks of the greater part as the
whole; and the praise was not theirs, ``for not even did this praise belong to
them, but to the grace of God''. To a church troubled by the delay he puts
Paul's meaning, ``why are ye troubled that Christ is not come? Nay, he is
come; and the Day is henceforth at the doors.'' At the word ``manifestation''
he notes, ```Revelation' is his word; implying that although He be not seen,
yet He is, and is present even now, and then shall appear'', and so ``there is
need of patience''. On v.~8, ``He shall confirm'' and
``unreprovable'' mark the Corinthians ``as still wavering, and liable to
reproof'': they are not yet blameless, and Paul says that Christ will make them
so. On v.~9 he gives Paul's reason: ``\,`God is faithful,' i.e. `true.' Now if
`true,' what things He hath promised He will also perform.'' Paul puts this
first, Chrysostom says, ``lest after the vehemence of the reproofs they might
fall into despair''; and where the call fails, ``unbelief cometh not of Him
that calleth, but of those who start away from Him'' (\work{Homilies on First
Corinthians} 2).

Aquinas's lecture on the same verses goes verse by verse. Those who wait for
the revelation are blessed now in hope. ``Unto the end'' is the end of each
one's life, and ``without crime'' means without mortal sin: \latin{qui sine
crimine invenitur in die mortis, sine crimine perveniet ad diem iudicii},
``whoever is found without crime on the day of death will come without crime
to the day of judgment''; Augustine, on the Gospel's command to all, ties each
one's death to the judgment in the same way. And God would not seem faithful ``if he called us to
the fellowship of his Son and denied us, as far as lies in him, the things by
which we could reach him''; the fellowship is had \latin{in praesenti per
similitudinem gratiae}, ``in the present by the likeness of grace'', and
\latin{in futuro per participationem gloriae}, ``in the future by sharing in
glory'' (\work{Super I ad Corinthios} c.~1, lect.~1).

\subsection{``The Lord will give goodness'': the earth and its fruit}

The Communion Antiphon sings Psalm 85 (84):13, ``the Lord will give goodness:
and our earth shall yield her fruit''. The first interpretation reads the
verse with Bellarmine, the third with Augustine. Bellarmine, who paraphrases
the Alleluia verse's petition as ``Send us your Son for a Savior'', reads it
of the Incarnation: the Lord sends his Spirit to overshadow a virgin, and
``our land, which was never ploughed nor sown, and was altogether an untouched
virgin, will yield her fruit'' (Ps 84, expl.~7 and~12). Augustine, expounding
the same verse, reads it after ``Truth is sprung out of the earth'', which in
this place he takes as confession: the sinner who confesses makes truth spring
from the earth. Then, \latin{Tua ergo peccata potes habere; fructum bonum non habebis, nisi ille dederit cui
confiteris}, ``Your sins, then, you can have of your own; good fruit you will
not have unless he gives it to whom you confess''. His psalter read
\latin{suavitatem}, ``sweetness'', and the sweetness the Lord gives is ``the
sweetness of working righteousness, so that righteousness shall begin to
delight thee, whom before unrighteousness delighted'' (84.15). Augustine had
made the same point at the Alleluia verse: the one to whom God shows mercy ``cannot indulge in pride'', since by showing him his salvation
God ``persuadeth him that whatever good man has, he hath not but from Him''
(84.9).

\subsection{The Missal's prayers: a will, a watch and gifts}

The third interpretation takes from the Entrance Antiphon the trust of its
second verse, ``In thee, O my God, I put my trust; let me
not be ashamed'', with Augustine, who hears the Church say: ``O my God, from
trusting in myself I was brought even to this weakness of the flesh'', and
``now, therefore, `in Thee I trust, I shall not be ashamed''' (24.2).

The Collect's petition is read three ways. In the first interpretation it
asks for the will to go out to meet a Christ who is still coming. In the
second its just works are close to the second of Gregory's marks of watching,
keeping by works what one believes. In the third it is a prayer for the will itself, \latin{Da \ldots\
hanc tuis fidelibus voluntatem}: the works follow from a will that is asked of
God. The first interpretation adds Aquinas on Isaiah's ``Thou hast met him
that rejoiceth'', whose verb the Collect shares. God, he says, goes out to meet in help the
one who is still asking or still on the way to him; and in a mystical sense
he names Simeon, the just and devout man of Jerusalem ``waiting for the
consolation of Israel'' (Lk 2:25; \work{Expositio super Isaiam} c.~64).

The Prayer over the Offerings offers gifts \latin{de tuis \ldots\ collata
beneficiis}, from what God has bestowed, as the servants in Bede's comment
work by the grace the master gave them. Its second petition, that what
devotion does in this passing time become the reward of eternal redemption,
is in the second interpretation a prayer of those who keep the watch with
works.

At the end of the Preface those who give thanks call themselves
\latin{vigilantes}, those who keep watch, from the verb of the Gospel's
\latin{vigilate}; what they now dare to wait for they call \latin{promissum},
what was promised, to be received openly when Christ comes in glory.

The Prayer after Communion says that by the mysteries God already teaches
those who live among passing things to love what is heavenly and to cling to
what lasts. Gregory, having said that death is certain for all, tells his
hearers not to let the care of earthly things weigh them down: \latin{Nolite
ergo attendere quid habetis, sed quid estis}, ``do not consider what you have,
but what you are'' (13.6). Gregory exhorts his hearers; the prayer says that
God teaches. The solemn blessing, where it is used, asks God to make the
faithful firm in faith, glad in hope and active in charity.

\subsection{Three answers, four senses}

All three interpretations take the first coming as given and the second as
awaited. The first asks about the comings themselves; the second and third
ask about the time between them, the one what it requires of those who wait,
the other what they can rely on while they wait. Two disagreements are real
and stand: Jerome and Aquinas on the literal sense of \latin{Descendisti}, and
Augustine and Bede on where the weight of the command falls. Jerome's patience
and Aquinas's permission are two compatible accounts of why God ``made us to
err''.

\begin{comparisontable}{}{Rend the heavens}{The watch at the door}{The clay and the potter}
Literal & A people asks God to come down as he came down to free Israel from Egypt (Aquinas);
the psalmist asks the Shepherd to come and save the vine & Jesus tells four
disciples, and all through them, to stay awake for an hour withheld so that
they live each day as if judged the next (Jerome) & A people confesses its
sin and appeals to God as Father and maker; Paul rests the Corinthians'
perseverance on God's faithfulness\\
Allegorical & The descent is the Word made flesh (Jerome); the son of man is
Christ (Augustine, Bellarmine); the first coming was unseen, like rain on a
fleece (Cyril) & The man gone abroad is Christ ascended, still present to his
Church; the servants' authority is the Spirit's grace; the doorkeeper the
pastors (Bede) & The face God is asked to show was veiled with ``a cloud of
flesh'' and acknowledged when Christ rose (Augustine 79.3); the called share
the Son's fellowship now by grace (Aquinas)\\
Moral & The believer prays the cry as his own, asks to be turned, and asks
for the will to go out to meet Christ & Each guards the door of his heart and
gives account for himself (Bede), keeps by works what he believes, and does
not despair of a later watch (Gregory) & The one shown mercy cannot be proud;
his sins are his own and his good fruit God's gift (Augustine)\\
Anagogical & The second coming is open and in glory (Cyril); Christ given as
man will be given as God (Augustine) & The world's last day takes each as his
own last day finds him (Augustine); found without grave sin at death, he comes
so to judgment (Aquinas) & Christ confirms to the end those who wait, and
fellowship by grace becomes a sharing in glory (Aquinas)\\
\end{comparisontable}
